% Ad Atticum 1.16 --- headnote
% ad-atticum-01-16, July 61 BC, written at Rome

\textit{Cicero to Atticus, written at Rome in July 61 BC ---
the long letter on the Clodius trial verdict and its
political aftermath. The Bona Dea bill of the previous
spring, sabotaged by the consul Piso (Att 1.13--14), had
been rescued by Hortensius's intervention only to be carried
in a form that delivered the case to a panel of judges
plainly buyable. Clodius was acquitted on the Ides of May,
and the present letter is Cicero's account of how it
happened: \textsection 1--2, the trial procedure and
Hortensius's miscalculation; \textsection 3--5, the bench
itself --- ``never was there a fouler bench in a gambling-
school: spotted senators, naked equites, tribunes not so
much `of the treasury' as `of the public account', i.e.
beggars''; the heroic minority of twenty-five judges who
voted to convict and the thirty-one ``whom hunger more than
fame moved''; the great moment when, as Cicero rose to
testify, the judges stood and offered their throats to
Clodius for Cicero's life. \textsection 6--9 is the political
diagnosis: the senate's authority is wounded but not
broken; ``twice was Lentulus acquitted, twice Catiline; this
man, the third now, has been let loose by the judges upon
the commonwealth.''}

\textit{The famous heart of the letter is \textsection 10
--- the verbatim altercation in the senate between Cicero
and Clodius. Six exchanges of riposte, each ending in
Cicero's punch line: ``You bought a house.'' / ``You would
think he was saying, `you bought judges'.'' --- ``They did
not believe you when you swore.'' / ``Twenty-five judges
believed me; thirty-one, since they took the money first,
did not believe you with anything.'' The closing
paragraphs (\textsection 11--13) report Cicero's standing
with the boni and with the urban plebs (where the rumour
that Pompey loves Cicero has made the bearded conspirators
call him ``Gnaeus Cicero'' in the streets); the bribery-
ridden consular elections of 60 BC just postponed by
Lurco's new bribery law; and the closing barb that, if
Pompey's man Afranius is elected, the consulship that
Curio called ``apotheosis'' will become a ``Bean Mime.''
The closing pieces of business --- a request for the
landscape and poetry of Atticus's Amaltheum at Buthrotum,
which Cicero plans to copy at Arpinum --- bring the letter
back to its private base.}
